% Fragment partage : inclus par compte-certifie-officiel.tex (dossier autonome)
% et par dossier-conception-alanya.tex (dossier client assemble).
% Ne pas y mettre de preambule ni de \part : c'est l'appelant qui les pose.

\chapter{Modèle de données}

\section{Ce que le socle a déjà posé}

\begin{constat}[Rien à recréer]
La migration~023 porte déjà sur \code{users}~: \code{verification\_status} (six états),
\code{verified\_until}, \code{verification\_reminder\_sent}, et l'index
\code{idx\_users\_verified\_until}. Le \code{verificationScheduler} de l'étape~1 assure déjà
la relance à J\,--\,30 et le passage en \emph{expiré} à l'échéance.
\end{constat}

Ce document n'ajoute que le \textbf{dossier} et ses pièces --- plus les tables commerciales,
laissées inactives.

\section{Migration \code{027\_verification.sql}}

\begin{lstlisting}[language=SQL, caption={Dossiers et pièces}]
-- 027_verification.sql -- dossiers de verification
-- Requiert 023 (socle de compte).

CREATE TABLE IF NOT EXISTS verification_request (
  id           BIGINT   NOT NULL AUTO_INCREMENT,
  alanyaID     INT      NOT NULL,
  target_type  TINYINT  NOT NULL DEFAULT 1
                 COMMENT 'Genre vise : 0=personnel notable 1=business',
  claimed_name VARCHAR(160) NULL COMMENT 'Raison sociale telle que declaree',
  status       TINYINT  NOT NULL DEFAULT 0
                 COMMENT '0=en attente 1=piece demandee 2=approuve 3=refuse 4=annule',
  reviewed_by  INT      NULL,
  decided_at   DATETIME NULL,
  reason       VARCHAR(255) NULL COMMENT 'Obligatoire si status = 3',
  created_at   DATETIME NOT NULL DEFAULT CURRENT_TIMESTAMP,
  PRIMARY KEY (id),
  KEY idx_vr_file (status, created_at),
  KEY idx_vr_user (alanyaID, created_at),
  CONSTRAINT fk_vr_user     FOREIGN KEY (alanyaID)
    REFERENCES users(alanyaID) ON DELETE CASCADE,
  CONSTRAINT fk_vr_reviewer FOREIGN KEY (reviewed_by)
    REFERENCES users(alanyaID) ON DELETE SET NULL
) ENGINE=InnoDB DEFAULT CHARSET=utf8mb4 COLLATE=utf8mb4_unicode_ci;

CREATE TABLE IF NOT EXISTS verification_document (
  id           BIGINT       NOT NULL AUTO_INCREMENT,
  request_id   BIGINT       NOT NULL,
  doc_type     TINYINT      NOT NULL
                 COMMENT '0=registre commerce 1=piece identite 2=adresse 3=notoriete',
  storage_key  VARCHAR(255) NOT NULL COMMENT 'Chemin du fichier chiffre',
  mime         VARCHAR(80)  NOT NULL,
  size_bytes   BIGINT       NOT NULL,
  -- Echeance portee par la ligne : plus simple qu'un script de retention
  -- separe, et lisible par quiconque ouvre la table.
  purge_after  DATETIME     NULL COMMENT 'Renseigne a la decision : +90 jours',
  purged_at    DATETIME     NULL,
  uploaded_at  DATETIME     NOT NULL DEFAULT CURRENT_TIMESTAMP,
  PRIMARY KEY (id),
  KEY idx_vd_purge (purge_after, purged_at),
  CONSTRAINT fk_vd_request FOREIGN KEY (request_id)
    REFERENCES verification_request(id) ON DELETE CASCADE
) ENGINE=InnoDB DEFAULT CHARSET=utf8mb4;

-- Journal d'acces : chaque ouverture d'une piece est tracee.
CREATE TABLE IF NOT EXISTS verification_document_access (
  id          BIGINT   NOT NULL AUTO_INCREMENT,
  document_id BIGINT   NOT NULL,
  admin_id    INT      NOT NULL,
  accessed_at DATETIME NOT NULL DEFAULT CURRENT_TIMESTAMP,
  ip          VARCHAR(45) NULL,
  PRIMARY KEY (id),
  KEY idx_vda_doc (document_id, accessed_at),
  CONSTRAINT fk_vda_doc   FOREIGN KEY (document_id)
    REFERENCES verification_document(id) ON DELETE CASCADE,
  CONSTRAINT fk_vda_admin FOREIGN KEY (admin_id)
    REFERENCES users(alanyaID)
) ENGINE=InnoDB DEFAULT CHARSET=utf8mb4;
\end{lstlisting}

\begin{note}[Pourquoi \code{purge\_after} sur la ligne]
Porter l'échéance dans l'enregistrement lui-même évite un script de rétention séparé, qui
serait la première chose oubliée lors d'un déploiement. Le même ordonnanceur que
l'expiration balaie \code{purge\_after <= NOW() AND purged\_at IS NULL}.
\end{note}

\begin{alerte}[Le journal d'accès n'est pas une option]
Des pièces d'identité consultables par n'importe quel administrateur sans trace, c'est le
genre de détail qui ne se remarque qu'après un incident. La table pèse peu et se purge avec
le reste.
\end{alerte}

\section{La machine à états du dossier}

\begin{center}
\begin{tikzpicture}[
  x=1cm, y=1cm,
  e/.style={draw=AlanyaLine, fill=AlanyaSurface, rounded corners=3pt,
            font=\sffamily\small, inner sep=6pt, minimum height=0.85cm, text=AlanyaInk},
  ok/.style={e, draw=AlanyaSuccess, line width=0.9pt, fill=AlanyaSuccess!8},
  ko/.style={e, draw=danger, fill=danger!6},
  fin/.style={e, draw=AlanyaIndigo, line width=0.9pt, fill=AlanyaContainer,
              text=AlanyaIndigoDark},
  a/.style={-{Stealth[length=2.2mm]}, draw=AlanyaGray, semithick},
  l/.style={font=\sffamily\scriptsize, text=AlanyaGray, fill=white, inner sep=2pt}
]
  \node[e]   (att) at (0,0)      {En attente};
  \node[ok]  (app) at (4.8,0)    {Approuvé};
  \node[fin] (ver) at (10.2,0)   {\code{verification\_status = 2}};
  \node[e]   (pie) at (0,-2.5)   {Pièce demandée};
  \node[ko]  (ref) at (4.8,-2.5) {Refusé};
  \node[ko]  (ann) at (10.2,-2.5){Annulé};

  \draw[a] (att) -- node[l, above] {approuve}        (app);
  \draw[a] (app) -- node[l, above] {pose l'échéance} (ver);

  % Aller-retour vers « pièce demandée » : deux traits parallèles, pour que
  % le retour ne se confonde pas avec l'aller.
  \draw[a] ([xshift=-9pt]att.south) -- node[l, left, pos=0.5] {pièce manquante}
           ([xshift=-9pt]pie.north);
  \draw[a] ([xshift=9pt]pie.north)  -- node[l, right, pos=0.5] {complété}
           ([xshift=9pt]att.south);

  \draw[a] (att.south east) -- node[l, pos=0.55, above right, inner sep=1pt]
           {refusé + motif} (ref.north west);

  % Abandon : route par le bas, sous les autres etats, pour ne rien croiser.
  \draw[a] (att.west) -- (-1.9,0) -- (-1.9,-4.1)
           -- node[l, below, pos=0.62] {abandon par le demandeur} (10.2,-4.1)
           -- (ann.south);
\end{tikzpicture}
\end{center}

\begin{note}[Un seul dossier ouvert à la fois]
Un compte ne peut avoir qu'une demande en \emph{attente} ou \emph{pièce demandée}. Sans cette
règle, un demandeur impatient en dépose trois et trois administrateurs les instruisent en
parallèle. Contrainte applicative --- une clé unique partielle n'existe pas en MySQL.
\end{note}

\chapter{Sécurité des pièces}

\section{Chiffrement au repos}

Les pièces sont des \textbf{données personnelles sensibles}. Trois exigences~:

\begin{center}
\begin{tabularx}{\linewidth}{@{}p{4.6cm}X@{}}
\toprule
\thd{Chiffrées au repos} & Chiffrement par fichier, clé hors de la base. Un accès en lecture
au disque ne doit pas suffire. \\
\thd{Jamais servies en statique} & \code{server.js} sert \code{/uploads} en statique
(\code{express.static}). Les pièces ne doivent \textbf{pas} y vivre~: elles passent par un
point d'entrée authentifié qui déchiffre à la volée. \\
\thd{Journalisées} & Chaque ouverture écrit une ligne avec l'administrateur et l'horodatage. \\
\bottomrule
\end{tabularx}
\end{center}

\begin{alerte}[Le piège du dossier \code{uploads}]
\code{src/middleware/upload.js} range tout sous \code{uploads/} et \code{server.js} expose
ce dossier en statique. Déposer une pièce d'identité par le chemin habituel la rendrait
\textbf{publiquement téléchargeable} pour qui connaît l'URL. Il faut un stockage distinct,
hors de l'arborescence servie.
\end{alerte}

\section{La purge}

Un balayage quotidien, sur le patron du \code{verificationScheduler}~:

\begin{lstlisting}[language=SQL, numbers=none]
SELECT id, storage_key FROM verification_document
WHERE purge_after IS NOT NULL
  AND purge_after <= NOW()
  AND purged_at IS NULL
LIMIT 200;
\end{lstlisting}

Le fichier est effacé du stockage, puis \code{purged\_at} est renseigné. La ligne survit
--- elle atteste qu'une pièce a existé et qu'elle a été détruite, ce qui est précisément ce
qu'on veut pouvoir montrer.

\chapter{Abonnement et paiement --- conçus, dormants}

\section{Pourquoi les concevoir maintenant}

Parce que la version~1 est gratuite, ces tables ne servent à rien aujourd'hui. On les décrit
tout de même, pour deux raisons~: leur forme conditionne des choix qu'on fait dès à présent
(le déblocage de fonctionnalités), et les ajouter plus tard sur une base peuplée coûte plus
cher que de les prévoir vides.

\begin{lstlisting}[language=SQL, caption={\code{027\_verification.sql} --- tables inactives en v1}]
CREATE TABLE IF NOT EXISTS subscription (
  id                 BIGINT   NOT NULL AUTO_INCREMENT,
  alanyaID           INT      NOT NULL,
  plan               TINYINT  NOT NULL DEFAULT 0,
  status             TINYINT  NOT NULL DEFAULT 0
                       COMMENT '0=inactif 1=actif 2=en retard 3=resilie',
  current_period_end DATETIME NULL,
  reminder_sent      TINYINT  NOT NULL DEFAULT 0,
  created_at         DATETIME NOT NULL DEFAULT CURRENT_TIMESTAMP,
  PRIMARY KEY (id),
  KEY idx_sub_user (alanyaID, status),
  KEY idx_sub_echeance (current_period_end),
  CONSTRAINT fk_sub_user FOREIGN KEY (alanyaID)
    REFERENCES users(alanyaID) ON DELETE CASCADE
) ENGINE=InnoDB DEFAULT CHARSET=utf8mb4;

CREATE TABLE IF NOT EXISTS payment_transaction (
  id              BIGINT       NOT NULL AUTO_INCREMENT,
  subscription_id BIGINT       NOT NULL,
  provider_ref    VARCHAR(120) NULL,
  -- Le fournisseur rejouera ses webhooks. Sans cette contrainte unique,
  -- double credit garanti.
  idempotency_key VARCHAR(120) NOT NULL,
  amount_minor    BIGINT       NOT NULL,
  currency        CHAR(3)      NOT NULL DEFAULT 'XAF',
  status          TINYINT      NOT NULL DEFAULT 0,
  created_at      DATETIME     NOT NULL DEFAULT CURRENT_TIMESTAMP,
  PRIMARY KEY (id),
  UNIQUE KEY uq_pt_idempotence (idempotency_key),
  KEY idx_pt_sub (subscription_id, created_at),
  CONSTRAINT fk_pt_sub FOREIGN KEY (subscription_id)
    REFERENCES subscription(id) ON DELETE CASCADE
) ENGINE=InnoDB DEFAULT CHARSET=utf8mb4;
\end{lstlisting}

\section{Renouvellement manuel, pas prélèvement}

\code{subscription} est modélisée pour un renouvellement \textbf{manuel} avec relances, et
non pour un prélèvement automatique. Sur mobile money, les autorisations expirent et les
échecs sont silencieux~; une formule annuelle relancée est plus réaliste qu'un auto-débit
qui casse sans prévenir. D'où \code{reminder\_sent}, exactement comme sur la vérification.

\section{La règle à ne jamais enfreindre}

\begin{alerte}[Ce qui débloque une fonctionnalité, c'est l'abonnement]
Aucune fonctionnalité ne doit être conditionnée à \code{verification\_status}. Les lieux
multiples de l'étape~4, par exemple, se débloquent par \code{subscription.status = 1}.

Si le badge ouvrait des portes, il deviendrait un produit qui se vend, et le retirer pour
défaut de paiement reviendrait à retirer une attestation d'identité à quelqu'un qui est
toujours la même personne. C'est le recouplage que le §4 du document d'amorce écarte, et
c'est ce qui permet aujourd'hui de livrer la certification gratuitement sans rien fermer.
\end{alerte}

\section{Une contrainte de plateforme à anticiper}

Un badge payant est un service numérique~: l'App~Store et le Play~Store peuvent exiger un
achat in-app, avec leur commission. C'est la raison pour laquelle les offres équivalentes
coûtent plus cher sur mobile que sur le web. À trancher \textbf{avant} la soumission, pas
pendant --- et c'est un argument de plus pour que la v1 reste gratuite.

\chapter{Impact fichier par fichier}

\section{Backend}

\begin{center}
\begin{tabularx}{\linewidth}{@{}p{6.2cm}p{1.6cm}X@{}}
\toprule
\thd{Fichier} & \thd{Nature} & \thd{Objet} \\
\midrule
\code{migrations/027\_verification.sql} & nouveau & Dossiers, pièces, journal, tables dormantes \\
\code{src/services/documentVault.js} & nouveau & Chiffrement, déchiffrement à la volée, purge \\
\code{src/controllers/verificationController.js} & nouveau & Dépôt, suivi, annulation \\
\code{src/controllers/admin/verification.js} & nouveau & File, instruction, décision \\
\code{src/routes/verification.js} & nouveau & \code{/api/verification} \\
\code{src/routes/admin.js} & modifié & File et décision \\
\code{src/services/verificationScheduler.js} & modifié & Ajout du balayage de purge \\
\code{server.js} & modifié & Montage des routes \\
\bottomrule
\end{tabularx}
\end{center}

\section{Administration et mobile}

\begin{center}
\begin{tabularx}{\linewidth}{@{}p{6.2cm}X@{}}
\toprule
\thd{Fichier} & \thd{Objet} \\
\midrule
\code{app/(dashboard)/verifications/page.tsx} \emph{(nouveau)} & File des dossiers \\
\code{app/(dashboard)/verifications/[id]/page.tsx} \emph{(nouveau)} & Instruction \\
\code{app/(dashboard)/layout.tsx} & Entrée de navigation, avec le compteur en attente \\
\code{lib/screens/profile/verification\_screen.dart} \emph{(nouveau)} & Dépôt et suivi \\
\code{lib/screens/profile/settings\_screen.dart} & Entrée « Faire vérifier mon compte » \\
\bottomrule
\end{tabularx}
\end{center}

\chapter{Points de vigilance et questions ouvertes}

\section{Chantiers nommés}

\begin{enumerate}[leftmargin=1.6em]
  \item \textbf{Le coffre à documents} --- chiffrement, service authentifié, purge, journal.
        Rien de cela n'existe~; c'est la brique la plus sensible des cinq étapes.
  \item \textbf{La sortie de \code{uploads/}} --- les pièces ne doivent pas transiter par le
        chemin de téléversement habituel, servi en statique.
\end{enumerate}

\section{Questions ouvertes}

\begin{enumerate}[leftmargin=1.6em]
  \item \textbf{Durée de validité par défaut}~: un an est l'usage, mais ce n'est pas tranché.
        Elle conditionne le rythme des relances et la charge d'instruction.
  \item La liste des pièces ci-dessus est une proposition~: elle demande une validation
        métier, notamment sur le RCCM comme pièce obligatoire.
  \item Un dossier refusé peut-il être redéposé immédiatement, ou après un délai~? Sans
        délai, un demandeur peut saturer la file.
  \item Que se passe-t-il si un compte vérifié \textbf{change de nom}~? L'attestation portait
        sur une identité~; un changement devrait au minimum alerter.
  \item Qui traite une \textbf{contestation} après refus~? Aucun circuit n'existe, et c'est
        précisément à cela que servent les 90~jours de rétention.
\end{enumerate}
